%! Mode:: "TeX:UTF-8"
%! TEX program = xelatex
\documentclass[zihao=-4,a4paper]{ctexart}
\usepackage[a4paper,top=2.5cm,bottom=2.5cm,left=2.6cm,right=2.6cm]{geometry}
\usepackage{booktabs}
\usepackage{float}
\usepackage{array}
\usepackage{fancyhdr}
\usepackage{lastpage}

% 页眉页脚：标注竞赛年份与题号，页码按“第 X 页 共 Y 页”格式
\setlength{\headheight}{12.65pt}
\pagestyle{fancy}
\fancyhf{}
\fancyhead[C]{\small 2026 年全国大学生数学建模竞赛 C 题}
\fancyfoot[C]{\small AI 工具使用详情 · 第 \thepage\ 页\quad 共 \pageref{LastPage} 页}
\renewcommand{\headrulewidth}{0.4pt}
\renewcommand{\footrulewidth}{0pt}

\setlength{\parindent}{2em}
\setlength{\parskip}{2pt}

\begin{document}

\begin{center}
  {\Large\bfseries AI 工具使用详情}\\[4pt]
  {\small 2026 年 9 月 13 日}
\end{center}

\section*{一、使用工具与范围}

竞赛期间使用 OpenAI Codex（GPT-5 系列），以及 Claude Code 与 DeepSeek 组合（内置模型为
DeepSeek V4.1 Flash），主要用于题意梳理、程序辅助编写与调试、模型检查、结果整理和论文表达优化。
使用时间为 2026 年 9 月 12 日至 13 日。

提问方式以给出明确问题背景、数据口径与验证要求的分步提示为主，要求 AI 提供可复算的推导、
代码与依据；采纳原则为“先验证、后采纳”，任何建议均须经人工检查或独立核算后方可进入模型、
程序和论文。

AI 工具未直接决定模型口径与最终结论。题目解释、模型假设、关键参数、
正式计算结果及提交文件均由参赛队员检查确认。

\section*{二、主要使用内容}

\begin{table}[H]
\centering
\small
\begin{tabular}{p{2.4cm} p{6.0cm} p{6.0cm}}
\toprule
环节 & AI 辅助内容 & 人工处理 \\
\midrule

题意与建模 &
整理题目条件，辅助检查时间槽、储能状态、合同购电和结算关系 &
参赛队员确定最终口径，并通过简单算例验证 \\

程序实现 &
辅助编写数据处理、优化求解、滚动回放和结果核验程序 &
程序经短窗口测试及独立核算后再用于正式计算 \\

结果分析 &
辅助整理费用、电量、预测误差和不同策略之间的比较 &
正式数值均从最终计算结果读取，并与结果文件核对 \\

论文整理 &
辅助修改章节结构、图表说明和文字表达 &
公式、数据和结论均由参赛队员最终复核 \\

\bottomrule
\end{tabular}
\end{table}

\section*{三、主要提示方式与使用过程}

项目采用“提出任务—AI 分析—人工判断—程序验证—结果回查”的交互方式，
AI 不直接一次性生成最终方案，而是在不同阶段辅助发现问题、提出候选方法并检查实现。
以下列出部分具有代表性的交互过程。

\subsection*{1.\ 建模思路检查}

\textbf{参赛队员：}
请检查当前 C 题建模方案，从题意完整性、逻辑一致性、实验可行性、
稳健性和创新性几个方面指出可能存在的问题，不要直接重写模型。

\textbf{AI 辅助：}
梳理四问之间的递进关系，重点检查时间尺度、储能状态连续性、
预测信息可得范围和不同购电费用的结算关系，并提出需要进一步验证的疑点。

\textbf{人工处理：}
参赛队员结合题目原文逐项判断，只保留有题面依据的建议；
对于存在歧义的时间边界和结算方式，通过最小算例和结果模板进一步验证。

\subsection*{2.\ 时间槽与跨日状态检查}

\textbf{参赛队员：}
检查 144 个十分钟槽位、0:00 与 24:00 的对应关系以及跨日 SOC 传递，
特别注意是否存在整体平移一个时间槽的问题。

\textbf{AI 辅助：}
指出附件时间标签、模型槽位与结果模板之间可能存在的错位，
并整理不同时间映射口径及其影响。

\textbf{人工处理：}
参赛队员最终确定统一时间方案，并通过单槽测试、跨午夜测试和工作簿回读，
确认模型输入、计算结果和模板位置一致。

\subsection*{3.\ 购电与费用结算检查}

\textbf{参赛队员：}
如果计划购电量大于实际需求，应如何处理？
请区分计划购电、实际交付、未使用合同电和紧急购电，
并检查费用是否存在重复计算。

\textbf{AI 辅助：}
提示不能将未使用合同电简单解释为弃光，
建议分别记录合同购电量、实际进入母线的电量、未使用合同电和紧急购电量；
同时通过简单的合同调整例子检查多次修改计划时可能出现的重复收费问题。

\textbf{人工处理：}
参赛队员采用独立费用账本，
逐槽核对正常购电费、调整现金流和紧急购电费，
并通过构造算例验证最终合同量不会被重复收费。

\subsection*{4.\ 预测因果性检查}

\textbf{参赛队员：}
所有预测和滚动优化只能使用决策时刻已经获得的数据。
请检查当前实现是否存在未来信息泄漏。

\textbf{AI 辅助：}
检查历史负荷、光伏和电价的读取时间，
指出多日预测视野中“目标时段”和“决策时刻”混淆可能造成的未来数据误用。

\textbf{人工处理：}
参赛队员统一规定所有预测输入满足
\[
\texttt{observed\_at} \leq \texttt{decision\_time},
\]
并重新运行因果性测试和相关正式结果。

\subsection*{5.\ 日内滚动调整}

\textbf{参赛队员：}
附件 3 在 0:00、6:00、12:00 和 18:00 发布新的光伏预测，
如何利用这些信息调整未执行计划，同时避免修改已经执行的时段？

\textbf{AI 辅助：}
建议采用滚动优化结构：
保留已执行时段，使用最新可获得预测重优化未来区间，
并将不同版本之间的合同变化单独记录。

\textbf{人工处理：}
参赛队员建立四时点顺序回放，
通过逐级加入 6:00、12:00 和 18:00 预报的消融实验，
以实际总费用和紧急购电量判断各发布时间的信息价值。

\subsection*{6.\ 波动电价处理}

\textbf{参赛队员：}
问题四中未来真实电价在制定计划时不可获得，
请检查是否应直接使用附件中的真实价格进行优化。

\textbf{AI 辅助：}
指出直接使用未来真实价格会造成信息泄漏，
建议将“预测价格用于制定计划”和“真实价格用于事后结算”分开处理，
并保留完美信息结果作为不可执行的理论下界。

\textbf{人工处理：}
参赛队员仅使用正式评价期以前的数据选择价格预测方法，
正式期内用预测价优化、真实价结算，
并由独立核算器重新计算全年费用。

\subsection*{7.\ 程序与结果核验}

\textbf{参赛队员：}
先完成小规模测试，再运行全年模型。
请检查物理约束、费用计算、时间连续性和结果文件是否一致。

\textbf{AI 辅助：}
辅助编写和检查测试程序，
整理需要验证的能量平衡、SOC 边界、跨日连续性、费用闭合和工作簿回读项目。

\textbf{人工处理：}
测试失败的实现不进入正式计算；
正式结果由独立核算程序逐槽复算，并重新打开官方结果文件进行核对。

\section*{四、结果核验}

为检查模型设定、程序实现和结果文件之间是否保持一致，本文在正式计算完成后对各问结果进行了进一步核验。

问题一采用显式充放电互斥的 MILP 与主线性规划模型进行对照，并逐槽检查母线能量平衡、储能状态递推、SOC 上下界、充放电功率限制以及终端储电量约束。两种模型在总费用和储能轨迹上保持一致，用于确认主模型不存在由同时充放电或约束遗漏造成的异常结果。

问题二至问题四采用按时间顺序回放的方式进行检查，重点核对预测数据是否只使用决策时刻以前已经获得的信息，避免未来数据进入计划过程；同时检查跨日 SOC 是否连续传递，各时段储能状态是否始终满足物理边界。对于问题三和问题四，还分别核对多次计划调整的版本顺序，以及预测价格和实际结算价格的使用范围。

费用部分单独按照逐槽记录重新计算。问题二核对正常合同购电和紧急购电费用；问题三进一步检查上调、下调和紧急购电费用是否按相邻计划版本正确结算；问题四则重新使用实际波动电价计算最终费用，避免直接采用优化器中的目标函数作为正式结果。

在结果文件生成后，再对论文中的主要费用、电量、SOC 和指定日期结果与正式工作簿进行逐项核对，检查时间槽、单位和汇总结果是否一致。只有完成上述检查且未发现约束违反、费用不闭合或结果文件不一致后，相关结果才用于论文正文和最终提交。

\end{document}
